\begin{abstract}
    The diffusion of power was a phenomenon first described by Susan Strange in the context of International Political Economy (IPE) in 1996. The phenomenon encompasses the concepts of relational and structural power, which are essential for understanding the IPE discipline. Although described in the 1990s, the diffusion of power did not address the cyber domain, nor explicitly the technological revolutions that culminated in the global computer network. In cyberspace, various actors operate to impose their authority over individuals and services. In this context, at the beginning of the 21st century, the low-latency anonymous network known as "The Onion Router" emerged, being part of the so-called "Dark Web" — a region of cyberspace characterized by active efforts to shield communication, privacy, and specific users. This research aims, through journalistic articles published between 2007 and 2017, to analyze the diffusion of power in the TOR network across three dimensions: authority, control, and results. 

    \vspace{1em}
    \textbf{Keywords:} Diffusion of Power. Dark Web. TOR Network. Structural Power. Relational Power. International Political Economy.
\end{abstract}